\section{Forward pass: embedding, historia y condicionamiento no Markov}

El diagrama de la arquitectura puede dar la impresión equivocada de que «basta con encadenar tres tipos de token». En realidad, cómo se combinan los embedding, qué hidden state se usan para decodificar, qué contiene la context window y si se permite que la historia influya en la acción de un mismo state son decisiones que cambian la función aprendida. Esta sección reconstruye, a partir de las slide y del Q\&A de la clase, un forward pass implementable.

\subsection{Entrada, salida y context window}

La entrada del modelo es una ventana local de la trayectoria con a lo sumo $K$ timestep; si se cuenta por token, contiene como máximo alrededor de $3K$ token de return/state/action. La salida es la secuencia de predicciones en las posiciones de action correspondientes:
\[
(\widehat{a}_1,\widehat{a}_2,\ldots,\widehat{a}_T)
=f_{\theta}(\widehat{R}_{1:T},s_{1:T},a_{1:T-1}).
\]
Aquí el lado derecho no incluye la action actual verdadera $a_T$; de lo contrario se produciría label leakage. En el entrenamiento, la implementación suele evitar la fuga mediante un shift o seleccionando hidden positions específicas.

\lecturefigure{slide-09-forward-pass.jpg}{Forward pass: entra la secuencia de return/state/action, sale la action predicha y la attention se calcula sobre un context de longitud $K$.}{Video oficial de Stanford Online, 14:56.}

\begin{importantbox}{Costo computacional de la context length}
En la self-attention estándar, el costo en tiempo y en matrices intermedias respecto del número de token es aproximadamente $O((3K)^2)$. Aumentar $K$ aporta una historia más larga, pero incrementa la memoria de GPU, la latencia y la dificultad de optimización. La ablation con $K=1$ que se presenta más adelante verifica si la historia contribuye realmente al desempeño, en lugar de tomar un modelo más grande como la respuesta.
\end{importantbox}

\subsection{Sumar embedding no equivale a concatenation}

Sean el state embedding, el token-type embedding y el timestep embedding $e_s(s_t)$, $e_{\mathrm{type}}$ y $e_t(t)$, respectivamente; una implementación posible es
\[
x_t^{(s)}=e_s(s_t)+e_{\mathrm{type}=s}+e_t(t).
\]
La suma exige que todos los vectores tengan la misma dimensión y equivale a una traslación dentro del mismo espacio de representación; la concatenation, en cambio, amplía la dimensión y deja que las capas posteriores aprendan cómo combinar. Ambas pueden expresar la información, pero difieren en número de parámetros, estructura geométrica e inductive bias.

\teachervoice{En el Q\&A de la clase se preguntó específicamente «¿suma o concatenación?». El ponente respondió que son dos decisiones distintas de codificación de la función; la implementación original usa addition y se observó que resultaba más adecuada. Las notas advierten: no hay que interpretar la suma de embedding como una pérdida de identidad, pues el token type y el timestep siguen entrando en la representación mediante vectores aprendibles.}

\subsection{El forward pass a nivel de código}

Las fórmulas anteriores muestran qué vectores se suman para formar un token, pero todavía no responden cómo se intercalan en la implementación las tres series temporales, cómo se garantiza que la action actual no se filtre ni qué hidden position debe leer el decoder. La slide original completa ese flujo de datos con un fragmento de código. El pseudocódigo equivalente que sigue desarrolla, en orden, embedding, stacking, causal attention y action head; solo expresa el mecanismo, no se ata a la API de ningún repository concreto y tampoco oculta el padding ni la selección de posiciones, los dos pasos donde es más fácil equivocarse.

\lecturefigure{slide-10-forward-pass-code.jpg}{Fragmento de código del forward pass de Decision Transformer mostrado en la clase.}{Video oficial de Stanford Online, 17:24.}

\begin{lstlisting}[caption={Pseudocódigo mínimo del forward pass de Decision Transformer},label={lst:dt-forward}]
def forward(returns_to_go, states, actions, timesteps, mask):
    time_emb = embed_time(timesteps)
    return_emb = embed_return(returns_to_go) + time_emb
    state_emb = embed_state(states) + time_emb
    action_emb = embed_action(actions) + time_emb

    tokens = interleave(return_emb, state_emb, action_emb)
    hidden = causal_transformer(tokens, attention_mask=mask)

    state_hidden = hidden[:, state_token_positions]
    action_pred = action_head(state_hidden)
    return action_pred
\end{lstlisting}

En el código, \texttt{interleave} intercala los token en el orden $(R,s,a)$; \texttt{state\_token\_positions} indica que el hidden state se toma en la posición donde ya se vieron el return y el state actuales, pero todavía no la action actual. Una implementación real también debe manejar distintas formas de state/action, padding mask, head continuo o discreto y sequence truncation.

\subsection{Por qué lo no Markov es una decisión deliberada}

La Markov property supone que el state actual ya contiene toda la información necesaria para predecir el futuro, de modo que la policy óptima puede depender solo de $s_t$. Sin embargo, muchas observaciones son en realidad una observation $o_t$ que solo ve una parte del estado oculto; la historia $o_{<t},a_{<t}$ ayuda a inferir la información oculta. Decision Transformer conserva explícitamente la history, lo que permite que «la misma observación, con distinta procedencia» corresponda a actions distintas.

\begin{knowledgebox}{Partial observability (observabilidad parcial)}
Si el estado real del entorno es $z_t$ y el modelo solo puede ver $o_t=g(z_t)$, entonces distintos $z_t$ pueden corresponder a la misma $o_t$. La trayectoria histórica puede servir como aproximación del belief-state: el modelo infiere el estado oculto actual a partir de las observaciones y acciones previas. En ese caso, el non-Markov conditioning no viola la teoría, sino que completa la información cuando la observation es insuficiente.
\end{knowledgebox}

\teachervoice{Un estudiante preguntó: si un mismo state recibe distintos embedding de tiempo en distintos timestep, ¿no se rompe la propiedad de Markov? El ponente respondió con claridad: «sí, aquí es deliberadamente no Markov», porque la action prediction y la partial observability pueden requerir la historia. Si la tarea es de verdad completamente Markov, la context ablation debería ayudar a determinar si la historia es solo redundante.}

\begin{warningbox}{La history también puede memorizar correlaciones erróneas}
Un context más largo puede recuperar el estado oculto, pero también puede hacer que el modelo dependa de patrones fortuitos del proceso de recolección de datos, como el ritmo de cierta behavior policy, la longitud de los episode o posiciones temporales específicas. Es necesario verificar, con pruebas entre policy distintas, entre estados iniciales distintos y con context ablation, que lo que aporta la historia sea información transferible y no una huella de los registros.
\end{warningbox}

\subsection{Resumen de la sección}

Lo esencial del forward pass no es solo el bloque Transformer, sino el intercalado de token, la combinación de embedding, el label shift, la context mask y el action head. La non-Markov history puede ser una ventaja en tareas parcialmente observables, pero una historia más larga también aumenta el costo computacional y el riesgo de spurious correlation.
